\section{GPT-3: Decoder-Only e In-Context Interface}

La parte de aplicaciones usa primero GPT-3 para mostrar la vía decoder-only. El modelo apila causal decoder blocks y se preentrena con next-token prediction sobre texto a gran escala. El objetivo de entrenamiento es simple, pero unifica muchas tareas como «dado un prefijo, predecir una continuación razonable». La instruction, los example y el context del prompt pasan a formar un mismo token stream.

\subsection{Autoregressive preentrenamiento y prompting}

Dada una secuencia de tokens $x_1,\ldots,x_T$, el modelo de lenguaje maximiza
\[
\sum_{t=1}^{T}\log p_\theta(x_t\mid x_{<t}).
\]
La causal mask garantiza que la posición $t$ solo vea el prefijo. \term{in-context learning} (aprendizaje en contexto) se refiere a que el modelo cambia su comportamiento en la tarea actual a partir de las instrucciones o los ejemplos dentro del prompt, sin actualizar sus parámetros. No es gradient-based training tradicional ni garantiza obtener reglas estables; el prompt format, el example order y el distribution shift influyen en el resultado.

\lecturefigure{slide-16-gpt3.jpg}{GPT-3 forma una interfaz de generación general con preentrenamiento causal decoder-only y prompt examples.}{Video oficial 18:30; Brown et al. 2020.}

\begin{knowledgebox}{Lectura de la figura: distinguir preentrenamiento, fine-tuning e in-context learning}
El preentrenamiento actualiza los parámetros para que el modelo aprenda una token distribution amplia; el fine-tuning sigue actualizando los parámetros con datos de la tarea; el in-context learning solo cambia el contexto de entrada y los parámetros permanecen iguales. Los ejemplos few-shot de la figura muestran que la task queda codificada en el prompt, no prueban que el modelo realmente ejecute un descenso de gradiente interno. Al evaluar, hay que comparar zero-shot, few-shot y el fine-tuned baseline, y controlar el token budget y la contamination.
\end{knowledgebox}

\teachervoice{El ponente usa GPT-3 para enfatizar «usar la red preentrenada como generator»: con una descripción de la tarea y unos cuantos example, produce una completion. La nota de clase es el cambio de interfaz: muchas tareas pasan de modificar el código del modelo a diseñar el contexto; esto aumenta la reutilización, pero también traslada los problemas de confiabilidad al prompt, a la evaluation y a los guardrails.}

\begin{warningbox}{El next-token objective no equivale a alineación con los hechos o con la intención}
Lo que optimiza el modelo de lenguaje es la likelihood condicional de los tokens, no «decir solo información verdadera» ni «cumplir el objetivo real del usuario». La escala y el prompting pueden mejorar el desempeño en las tareas, pero la exactitud factual, la calibration, la seguridad y la tool-use correctness requieren datos, entrenamiento y controles de sistema adicionales.
\end{warningbox}

\subsection{Resumen de la sección}

GPT-3 convierte el Transformer decoder-only en una prefix-to-continuation interface general. Muestra la escala del preentrenamiento y el in-context behavior, pero no elimina la definición de la tarea, la evaluación ni los límites de seguridad.
